\documentclass[conference]{IEEEtran}
\IEEEoverridecommandlockouts
%Template version as of 6/27/2024

\usepackage{cite}
\usepackage{amsmath,amssymb,amsfonts}
\usepackage{algorithmic}
\usepackage{graphicx}
\usepackage{textcomp}
\usepackage{xcolor}
\usepackage[utf8]{inputenc}
\usepackage[T1]{fontenc}
\usepackage{booktabs}

\def\BibTeX{{\rm B\kern-.05em{\sc i\kern-.025em b}\kern-.08em
    T\kern-.1667em\lower.7ex\hbox{E}\kern-.125emX}}

\begin{document}

\title{Detección Automática de Fracturas Óseas en Radiografías\\
Mediante YOLOv11 y RT-DETR con Preprocesamiento Fotométrico}

\author{
\IEEEauthorblockN{Lucia T. Capon Paul}
\IEEEauthorblockA{\textit{CEIA -- FIUBA} \\
Buenos Aires, Argentina}
\and
\IEEEauthorblockN{César Andrés Orellana}
\IEEEauthorblockA{\textit{CEIA -- FIUBA} \\
Buenos Aires, Argentina}
\and
\IEEEauthorblockN{Leandro Britez}
\IEEEauthorblockA{\textit{CEIA -- FIUBA} \\
Buenos Aires, Argentina}
}

\maketitle

% ─── ABSTRACT ────────────────────────────────────────────────────────────────
\begin{abstract}
La detección automatizada de fracturas óseas en imágenes radiográficas representa
un desafío clínico de alto impacto, donde la precisión y la velocidad de diagnóstico
son críticas. En este trabajo se presenta un prototipo basado en modelos de detección
de objetos preentrenados para la identificación y localización de fracturas en
radiografías de múltiples regiones anatómicas. Se comparan dos arquitecturas modernas:
YOLOv11 \cite{b3}, representativa del paradigma \textit{anchor-free} de una etapa,
y RT-DETR \cite{b4}, un detector basado en transformers de tiempo real. Ambos modelos
son ajustados (\textit{fine-tuned}) sobre el dataset público \textit{Bone Fracture
Detection} de Roboflow Universe \cite{b5}, que contiene 4.148 radiografías
distribuidas en seis clases anatómicas. Previamente al entrenamiento se aplica
un pipeline de preprocesamiento fotométrico que combina normalización robusta por
percentiles, ecualización adaptativa de histograma con límite de contraste (CLAHE)
y enmascaramiento no nítido (\textit{unsharp masking}), siguiendo la metodología
propuesta por Srinivasu \textit{et al.} \cite{b1}. La evaluación se realiza en
términos de mAP@0.5, mAP@0.5:0.95, F1-score y latencia de inferencia.
% [COMPLETAR CON RESULTADOS FINALES]
\end{abstract}

\begin{IEEEkeywords}
detección de objetos, fracturas óseas, YOLOv11, RT-DETR, CLAHE, \textit{unsharp
masking}, \textit{deep learning}, visión por computadora, radiografías.
\end{IEEEkeywords}

% ─── I. INTRODUCCIÓN ─────────────────────────────────────────────────────────
\section{Introducción}

La interpretación de imágenes radiográficas es una tarea clínica fundamental en el
diagnóstico de traumatismos musculoesqueléticos. El cuerpo humano posee 206 huesos
cuyas fracturas pueden presentarse como interrupciones transversales, oblicuas,
espirales o conminutas, muchas veces con escaso contraste respecto del tejido
circundante \cite{b1}. La identificación visual requiere experiencia especializada y
puede verse comprometida por factores como la carga de trabajo, la fatiga del
profesional o variaciones instrumentales en la adquisición de la imagen.

Los sistemas de visión por computadora basados en \textit{deep learning} han
demostrado capacidad para asistir en tareas de detección y localización de
patologías en imágenes médicas \cite{b2}. En particular, los detectores de objetos
permiten no solo determinar la presencia de una lesión sino también localizar su
posición mediante cuadros delimitadores (\textit{bounding boxes}), lo que
complementa el criterio clínico sin pretender reemplazarlo.

Un antecedente directo es el trabajo de Srinivasu \textit{et al.} \cite{b1}, que
evaluó YOLOv10 para detectar fracturas multiclase en radiografías con el mismo
dataset de Roboflow Universe. Los autores combinaron CLAHE y \textit{unsharp masking}
como preprocesamiento fotométrico y reportaron una mejora de mAP@0.5 de 0,948 a
0,961 respecto de imágenes sin procesar. Motivados por este antecedente, el presente
trabajo compara YOLOv11 y RT-DETR bajo un protocolo de partición fija y evalúa el
efecto del mismo pipeline de preprocesamiento sobre la capacidad de localización y
clasificación de fracturas.

El objetivo de este trabajo es desarrollar y evaluar un prototipo funcional capaz
de detectar automáticamente fracturas óseas en radiografías, comparando el
desempeño de YOLOv11 \cite{b3} y RT-DETR \cite{b4} en términos de métricas de
detección y latencia de inferencia. Los aportes principales son: (i) la aplicación
y justificación de un pipeline de preprocesamiento conservador adaptado a la
variabilidad del dataset, (ii) la comparación directa de un detector de una etapa
con un detector basado en transformers bajo el mismo protocolo experimental, y
(iii) el análisis por clase para identificar limitaciones del dataset.

% ─── II. METODOLOGÍA ─────────────────────────────────────────────────────────
\section{Metodología, Diseño y Desarrollo}

\subsection{Dataset}

Se utilizó el dataset \textit{Bone Fracture Detection} disponible públicamente en
Roboflow Universe (versión~1, licencia CC~BY~4.0) \cite{b5}. El dataset contiene
radiografías de múltiples regiones anatómicas con anotaciones en formato YOLO
(\textit{bounding boxes}). La Tabla~\ref{tab:dataset} resume la distribución de
imágenes y anotaciones por partición.

\begin{table}[htbp]
\caption{Distribución del dataset por partición}
\label{tab:dataset}
\begin{center}
\begin{tabular}{lccc}
\toprule
\textbf{Partición} & \textbf{Imágenes} & \textbf{Con bbox} & \textbf{Anotaciones} \\
\midrule
Train & 3.631 & 1.804 & 2.088 \\
Valid & 348   & 173   & 204   \\
Test  & 169   & 83    & 96    \\
\midrule
\textbf{Total} & \textbf{4.148} & \textbf{2.060} & \textbf{2.388} \\
\bottomrule
\end{tabular}
\end{center}
\end{table}

Las imágenes sin caja delimitadora se tratan como negativos verdaderos y permiten
al modelo aprender cuándo no debe activarse. No se hallaron duplicados exactos entre
particiones (verificado mediante hash SHA-256), lo que descarta fuga de datos
(\textit{data leakage}). La auditoría de anotaciones no detectó clases fuera de
rango ni geometrías inválidas.

El dataset original declara siete clases. Durante el análisis exploratorio se
comprobó que la clase \texttt{humerus} presentaba únicamente tres anotaciones en
entrenamiento y ninguna en validación ni test, constituyendo una anomalía residual
sin representatividad estadística. Esas anotaciones fueron unificadas
programáticamente con la clase \texttt{humerus fracture} por coherencia anatómica,
resultando en seis clases efectivas. La distribución de anotaciones en entrenamiento
se muestra en la Tabla~\ref{tab:clases}.

\begin{table}[htbp]
\caption{Distribución de anotaciones por clase en entrenamiento}
\label{tab:clases}
\begin{center}
\begin{tabular}{lcc}
\toprule
\textbf{Clase} & \textbf{Anotaciones} & \textbf{\%} \\
\midrule
fingers positive  & 531 & 25{,}43 \\
shoulder fracture & 360 & 17{,}24 \\
elbow positive    & 339 & 16{,}24 \\
forearm fracture  & 316 & 15{,}13 \\
humerus fracture  & 314 & 15{,}04 \\
wrist positive    & 228 & 10{,}92 \\
\midrule
\textbf{Total}    & \textbf{2.088} & \textbf{100} \\
\bottomrule
\end{tabular}
\end{center}
\end{table}

El dataset presenta un desbalance leve a moderado: la clase mayoritaria
(\texttt{fingers positive}) más que duplica a la minoritaria (\texttt{wrist
positive}). Esto se tiene en cuenta al analizar las métricas por clase, en
particular para determinar si el modelo logra detectar correctamente las
categorías con menor representación. El análisis exploratorio reveló además que
las imágenes tienen resoluciones heterogéneas (ancho medio: 412~px, alto medio:
486~px), niveles de brillo variados (media: 54/255) y que aproximadamente la mitad
de cada partición carece de fractura anotada, lo que refleja la prevalencia clínica
real.

\subsection{Preprocesamiento}

Las imágenes presentan variaciones de contraste y brillo atribuibles a diferencias
en el equipo, la posición del paciente y la exportación del archivo. Para mejorar
la visibilidad local de las estructuras óseas antes del entrenamiento, se aplicó un
pipeline de preprocesamiento determinista, igual sobre las tres particiones. No se
modificó la geometría de las imágenes ni las anotaciones.

El pipeline se evaluó comparando visualmente dos configuraciones. Se eligió la
configuración~P2 (realce conservador) sobre P1 (realce más agresivo) porque P1
amplificaba el granulado y los marcadores de los equipos, introduciendo artefactos
potencialmente confusos para el detector.

\subsubsection{Normalización robusta por percentiles}

Se escalan las intensidades al rango $[0, 255]$ descartando los valores extremos
(percentiles 1 y~99). Esto evita que píxeles anómalos muy oscuros o muy claros
dominen el histograma global, produciendo imágenes con rangos dinámicos comparables
entre radiografías de distintos equipos.

\subsubsection{CLAHE}

La ecualización adaptativa de histograma con límite de contraste (CLAHE) divide la
imagen en regiones de $8 \times 8$ píxeles y ecualiza el histograma local con un
límite de amplificación de 1{,}5. A diferencia de la ecualización global, el límite
de contraste evita la amplificación excesiva del ruido en regiones homogéneas,
mejorando la visibilidad de estructuras óseas de bajo contraste sin distorsionar
las zonas uniformes.

\subsubsection{Enmascaramiento no nítido (\textit{unsharp masking})}

El \textit{unsharp masking} refuerza las componentes de alta frecuencia (bordes)
restando a la imagen una versión suavizada de sí misma:

\begin{equation}
O(x,y) = I(x,y) + \gamma\bigl[I(x,y) - I_b(x,y)\bigr],
\label{eq:unsharp}
\end{equation}

donde $I(x,y)$ es la intensidad original, $I_b(x,y)$ la imagen suavizada con
$\sigma = 1{,}0$ y $\gamma = 0{,}25$ controla la magnitud del realce. Un valor
conservador de $\gamma$ resalta transiciones sutiles sin introducir bordes
artificiales. Los parámetros completos de la configuración~P2 se resumen en la
Tabla~\ref{tab:preproc}.

\begin{table}[htbp]
\caption{Parámetros del pipeline de preprocesamiento (configuración P2)}
\label{tab:preproc}
\begin{center}
\begin{tabular}{lc}
\toprule
\textbf{Parámetro} & \textbf{Valor} \\
\midrule
Normalización percentil bajo / alto & 1 / 99 \\
CLAHE clip limit                    & 1{,}5  \\
CLAHE tile grid size                & $8 \times 8$ \\
Unsharp $\sigma$                    & 1{,}0  \\
Unsharp $\gamma$                    & 0{,}25 \\
Calidad JPEG exportación            & 95     \\
\bottomrule
\end{tabular}
\end{center}
\end{table}

\subsection{Arquitecturas}

\textbf{YOLOv11} \cite{b3}: detector \textit{anchor-free} de una sola etapa de la
familia Ultralytics. Se utilizó la variante \texttt{yolo11m.pt} (medium, $\approx$20M
parámetros) preentrenada en COCO. YOLOv11 introduce mejoras en el bloque C3k2,
atención SPPF con cabezal desacoplado y una cadena de regresión de cuadros más
eficiente que sus predecesores.

\textbf{RT-DETR-L} \cite{b4}: detector basado en transformers diseñado para
inferencia en tiempo real. Se utilizó la variante \texttt{rtdetr-l.pt} (large,
$\approx$32M parámetros) preentrenada en COCO. RT-DETR emplea un \textit{backbone}
ResNet con un \textit{encoder} híbrido CNN-transformer y un \textit{decoder} de
atención cruzada que elimina la supresión de no máximos (NMS) del post-procesamiento.

La elección de estas dos arquitecturas responde a la necesidad de comparar el
paradigma de detección de una etapa con regresión directa (YOLOv11) frente al
paradigma basado en atención global (RT-DETR), analizando el compromiso entre
precisión de localización y latencia de inferencia.

\subsection{Configuración del Entrenamiento}

Ambos modelos se configuraron para un máximo de 100 épocas con tamaño de entrada
$640 \times 640$ píxeles. La Tabla~\ref{tab:hparams} detalla los hiperparámetros
de cada modelo. El ajuste fino (\textit{fine-tuning}) transfiere los pesos
preentrenados en COCO y reemplaza la cabeza de clasificación para adaptarla a las
seis clases del dataset.

\begin{table}[htbp]
\caption{Hiperparámetros de entrenamiento por modelo}
\label{tab:hparams}
\begin{center}
\begin{tabular}{lcc}
\toprule
\textbf{Parámetro}    & \textbf{YOLOv11}   & \textbf{RT-DETR-L}  \\
\midrule
Pesos base            & yolo11m.pt         & rtdetr-l.pt         \\
Épocas (máx.)         & 100                & 100                 \\
Early stopping        & --                 & patience = 20       \\
Optimizador           & auto (SGD)         & AdamW               \\
Batch size            & 16                 & 8                   \\
lr$_0$                & 0{,}01             & 0{,}0001            \\
lr$_f$                & 0{,}001            & 0{,}00001           \\
Weight decay          & 0{,}0005           & 0{,}0001            \\
Warmup epochs         & 3                  & 2                   \\
Mosaic aug.           & 1{,}0              & 0{,}0               \\
Mixup aug.            & 0{,}1              & 0{,}0               \\
Flip horizontal       & 0{,}5              & 0{,}5               \\
Flip vertical         & 0{,}0              & 0{,}0               \\
Conf. umbral inf.     & 0{,}25             & 0{,}25              \\
IoU NMS               & 0{,}70             & 0{,}70              \\
\bottomrule
\end{tabular}
\end{center}
\end{table}

Para RT-DETR se desactivaron \textit{mosaic} y \textit{mixup} por incompatibilidad
con la arquitectura basada en atención. Se empleó una tasa de aprendizaje
significativamente menor (0,0001 vs.\ 0,01) dado que los \textit{decoders} de tipo
DETR son sensibles a tasas elevadas durante el ajuste fino; el optimizador se fijó
explícitamente en AdamW, dado que el modo automático de selección de Ultralytics
(\texttt{optimizer=auto}) ignora la tasa de aprendizaje especificada por el usuario
y determina una propia. Asimismo, se detectó que el entrenamiento inicial de
RT-DETR divergía (\textit{loss} de clasificación con valores NaN) al utilizar
precisión mixta (AMP); este comportamiento es un problema conocido de los
\textit{decoders} basados en atención de tipo DETR bajo \texttt{fp16}, por lo que
se deshabilitó AMP para dicho modelo. Se incorporó además un criterio de
\textit{early stopping} (\texttt{patience=20}) para evitar cómputo innecesario una
vez alcanzada la meseta de convergencia; con este criterio, RT-DETR-L detuvo su
entrenamiento en la época 58, correspondiendo el mejor checkpoint a la época 38.
El entrenamiento se realizó en Google Colab con GPU NVIDIA A100 (80GB),
implementado con la librería Ultralytics. Los archivos de configuración completos
están disponibles en el repositorio (\texttt{configs/yolov11.yaml} y
\texttt{configs/model2.yaml}).

% ─── III. RESULTADOS ─────────────────────────────────────────────────────────
\section{Resultados Experimentales y Análisis}

Esta sección reporta los resultados obtenidos para RT-DETR-L sobre el conjunto de
test (169 imágenes, 96 anotaciones), evaluado con el mejor \textit{checkpoint}
(época 38) bajo el protocolo descrito en la Sección~II. La evaluación de YOLOv11m
sobre el mismo conjunto de test se encuentra en curso al momento de este informe;
la Tabla~\ref{tab:results} se completará con dichos valores en la versión final.

\begin{table}[htbp]
\caption{Comparación de métricas en el conjunto de test}
\label{tab:results}
\begin{center}
\begin{tabular}{lcccc}
\toprule
\textbf{Modelo} & \textbf{mAP@0.5} & \textbf{mAP@0.5:0.95} & \textbf{F1} & \textbf{Lat. (ms)} \\
\midrule
YOLOv11m  & -- & -- & -- & -- \\
RT-DETR-L & 0{,}2475 & 0{,}0802 & 0{,}2854 & 41{,}68 \\
\bottomrule
\end{tabular}
\end{center}
\end{table}

RT-DETR-L alcanzó una latencia mediana de 41{,}68~ms por imagen (P95:
46{,}38~ms; $\approx$24 FPS) sobre una GPU NVIDIA A100, sin \textit{batching}
(\textit{batch size}~=~1), lo que resulta compatible con un escenario de
asistencia diagnóstica no crítico en tiempo real, aunque debe evaluarse en
hardware de despliegue equivalente al clínico antes de extraer conclusiones
sobre viabilidad en producción.

\subsection{Análisis por Clase}

Dado que el mAP global promedia clases con soporte muy dispar (desde 6 hasta 43
instancias en test), se analiza el desempeño por clase en la
Tabla~\ref{tab:per-class}.

\begin{table}[htbp]
\caption{Métricas por clase de RT-DETR-L en el conjunto de test}
\label{tab:per-class}
\begin{center}
\begin{tabular}{lccc}
\toprule
\textbf{Clase} & \textbf{P} & \textbf{R} & \textbf{AP@0.5} \\
\midrule
forearm fracture  & 0{,}720 & 0{,}553 & 0{,}512 \\
humerus fracture  & 0{,}387 & 0{,}467 & 0{,}375 \\
shoulder fracture & 0{,}327 & 0{,}235 & 0{,}238 \\
wrist positive    & 0{,}093 & 0{,}167 & 0{,}176 \\
fingers positive  & 0{,}395 & 0{,}122 & 0{,}140 \\
elbow positive    & 0{,}000 & 0{,}000 & 0{,}043 \\
\midrule
\textbf{Media}    & \textbf{0{,}320} & \textbf{0{,}257} & \textbf{0{,}248} \\
\bottomrule
\end{tabular}
\end{center}
\end{table}

El modelo alcanza su mejor desempeño en \texttt{forearm fracture}
(AP@0.5~=~0,512) y \texttt{humerus fracture} (AP@0.5~=~0,375), clases con
anotaciones de mayor tamaño relativo y contorno óseo bien definido. En el
extremo opuesto, \texttt{elbow positive} obtiene Precision y Recall nulos al
umbral de confianza de inferencia (0,25): el modelo no logró aprender un
patrón discriminativo para esta clase, pese a contar con la tercera mayor
cantidad de anotaciones en entrenamiento (339), lo que sugiere que el
problema no es puramente de escasez de datos sino de variabilidad
morfológica o solapamiento visual con clases vecinas (a confirmar con la
matriz de confusión). \texttt{wrist positive}, pese a ser la clase con menor
soporte en test (6 instancias), no es la de peor desempeño, lo que indica que
el tamaño de muestra por sí solo no explica el patrón de errores observado.

\subsection{Dinámica de Entrenamiento}

RT-DETR-L mostró una convergencia rápida durante las primeras 30 a 38 épocas,
entrando luego en una meseta sin mejoras sostenidas en el \textit{fitness} de
validación, lo que activó el criterio de \textit{early stopping} en la época
58. Este comportamiento es consistente con la sensibilidad reportada de los
\textit{decoders} DETR a la capacidad del dataset: con 3.631 imágenes de
entrenamiento, el mecanismo de \textit{matching} húngaro dispone de señal
limitada para refinar las consultas (\textit{queries}) de las clases con
peor desempeño.

% TODO: agregar curvas de entrenamiento (Fig. training_curves.png), matriz de
%       confusión (confusion_matrix_normalized.png) y ejemplos cualitativos
%       (qualitative_predictions.png) como figuras una vez descargados desde
%       Drive al repositorio (carpeta results/).
% TODO: completar Tabla~\ref{tab:results} y esta sección con los resultados
%       de YOLOv11m para habilitar la comparación directa.

% ─── IV. CONCLUSIONES ────────────────────────────────────────────────────────
\section{Conclusiones}

% [TODO: completar con al menos 5 conclusiones basadas en los resultados]

\begin{enumerate}
    \item ...
    \item ...
    \item ...
    \item ...
    \item ...
\end{enumerate}

% ─── BIBLIOGRAFÍA ────────────────────────────────────────────────────────────
\begin{thebibliography}{00}

\bibitem{b1}
P.~N.~Srinivasu, G.~L.~A.~Kumari, S.~C.~Narahari, S.~Ahmed y A.~Alhumam,
``Exploring the impact of hyperparameter and data augmentation in YOLO V10 for
accurate bone fracture detection from X-ray images,''
\textit{Scientific Reports}, vol.~15, art.~9828, 2025.
doi: 10.1038/s41598-025-93505-4.

\bibitem{b2}
Z.~Su, X.~Adam, Y.~Zhao, H.~Fan, Z.~Dong y L.~Zhao,
``Skeletal fracture detection with deep learning: A comprehensive review,''
\textit{Diagnostics}, vol.~13, n.º~20, art.~3245, 2023.
doi: 10.3390/diagnostics13203245.

\bibitem{b3}
Ultralytics, ``YOLO11: State-of-the-Art Object Detection,'' 2024.
[En línea]. Disponible: https://github.com/ultralytics/ultralytics

\bibitem{b4}
Y.~Zhao, W.~Lv, S.~Xu, J.~Wei, G.~Wang, Q.~Dang, Y.~Liu y J.~Chen,
``DETRs Beat YOLOs on Real-time Object Detection,''
\textit{arXiv preprint arXiv:2304.08069}, 2023.

\bibitem{b5}
Roboflow Universe, ``Bone Fracture Detection Dataset,'' 2024.
[En línea]. Disponible: https://universe.roboflow.com/veda/bone-fracture-detection-daoon

\bibitem{b6}
R.~Y.~Ju y W.~Cai,
``Fracture detection in pediatric wrist trauma X-ray images using YOLOv8 algorithm,''
\textit{Scientific Reports}, vol.~13, art.~20077, 2023.
doi: 10.1038/s41598-023-47460-7.

\end{thebibliography}

\end{document}
